\subsection{Land Tenure Security}

As shown, the direct ecological impact of settlement schemes is driven by LULC change, happening for a variety of reasons.
The institutional and legal framework of property rights serves as a moderating mechanism.
Literature on land tenure security has shown that secure property rights are a prerequisite for investments into sustainable land management and agriculture including soil conservation \parencite{aliEnvironmentalGenderImpacts2014},
reduction of forest clearing \parencite{robinsonDoesSecureLand2014} and the planting of new trees \parencite{besleyLandReformPoverty2000, deiningerOverlappingLandRights2008}.
\cite{besleyLandReformPoverty2000} found that secure land tenure incentivizes long-term investments by reducing the risk of expropriation and creating access to credit markets needed for larger investments.
\cite{robinsonDoesSecureLand2014} also determined that the form of tenure, whether communal or private, is of little importance for ecological outcomes in Sub-Saharan Africa, further highlighting the importance of security.

The empirical evidence on this relationship between LTS and environmental outcomes in Sub-Saharan Africa is mixed and depends on the specific context \parencite{tsengInfluenceLandTenure2020,noufeImpactLandTenure2023}.
\cite{goldsteinProfitsPowerLand2008} found that in Ghana, individuals without political power face higher insecurity of land tenure, which leads to significantly shorter fallow periods.
Adding to this, \cite{deiningerOverlappingLandRights2008} showed that in Uganda, plots tilled by individuals with ownership rights planted significantly more trees than plots tilled by individuals only occupying the land.
In contrast, \cite{deiningerOverlappingLandRights2008} demonstrate that the presence of land titles alone does not guarantee positive outcomes in cases of overlapping land rights.
This coincides with \cite{naughton-trevesLandTenureTropical2014}, who found that de jure LTS is a basic condition for environmental conservation, with the actual quality and enforcement of rights determining the real outcomes.

In the specific case of Kenya and the post-colonial government throughout different forms of government, multiple administrative bodies and political actors have been involved in the allocation of land \parencite{ndunguReportComissionInquiry2004}.
The specific case of overlapping land rights can be seen in the illegal settlement schemes explained in the previous section,
with a report by the Kenyan government in 2004 documenting the same plot of land being allocated two or in some cases three times to different individuals \parencite{ndunguReportComissionInquiry2004}.
Furthermore, \cite{boonePromisedLandSettlement2021,booneLandConflictDistributive2012} found that many schemes were allocated without proper registration for years, creating insecurity through the political dependence on the state.
Adding to this, many of the plots allocated at the time were placed on already populated land, by displacing the original occupants \parencite{ndunguReportComissionInquiry2004},
and issuing land titles to the new occupants, often political allies, adding further to the insecurity of land tenure as found by \cite{booneLandConflictDistributive2012}.

The institutional disarray and de facto insecurity of tenure could negatively impact environmental efforts,
as the delay in title issuance and thereby a secure framework disincentivizes settlers from undertaking long-term investments like establishing a canopy of trees
or soil conservation measures, as shown by \cite{aberraDoesLegalRepresentation2021}.
Instead, a focus on short-term maximization of profits through logging or intensive agriculture is incentivized, which could lead to long-term degradation of land.
This leads to the second hypothesis:

Hypothesis 2 (H2): The negative impact of settlement schemes on local environmental conditions is moderated by the duration of insecurity experienced between land allocation and title issuance,
with lower security leading to more vegetation loss within program boundaries compared to neighboring non-program areas.